\documentclass[twocolumn]{aastex631}

\usepackage{amsmath}
\usepackage{natbib}
\usepackage{graphicx}
\usepackage{multirow}

\shorttitle{Calibration Audit of Transient Classifiers}
\shortauthors{Kailas}

\begin{document}

\title{Calibration Audit of Production Astronomical Transient Classifiers on ZTF Alert Streams}

\author{Pallavi Kailas}
\affiliation{Independent Researcher}
\email{pallavi.kailas@example.com}

\begin{abstract}
Production transient classifiers deployed on the Zwicky Transient Facility (ZTF) alert stream are routinely used to prioritize spectroscopic follow-up, yet their probability outputs have never been systematically validated as calibrated. We present the first calibration audit of four classifiers across three systems---ALeRCE's Balanced Random Forest, Fink's Random Forest and SuperNNova, and the NEEDLE framework---using spectroscopically confirmed transients from the ZTF Bright Transient Survey as ground truth. We compute Expected Calibration Error (ECE) with equal-mass binning, classwise reliability diagrams, and Brier scores, then apply temperature scaling with five-fold cross-validation. We find that ALeRCE exhibits moderate miscalibration (ECE $= 0.271$), which temperature scaling with $T = 0.357$ reduces to 0.097 (65\% improvement); applying a recalibrated 0.8 probability threshold increases usable candidates from 20 to 425 objects while maintaining 91\% precision, a 21-fold operational gain. Fink implements implicit selective classification: the Random Forest abstains on 94\% of objects via a minimum-epoch eligibility gate, achieving conditional ECE $= 0.263$ [0.211, 0.409] on the 6\% where predictions are made; SuperNNova abstains on 36\% of objects, with conditional ECE $= 0.183$ [0.155, 0.220] on predictions, which temperature scaling improves to 0.118. NEEDLE appears well-calibrated in aggregate (ECE $= 0.073$) but per-class analysis reveals class-asymmetric miscalibration: SLSN-I predictions are overconfident while SN predictions are underconfident, driven by inverse-frequency class weighting. Global temperature scaling worsens NEEDLE's calibration, demonstrating that a single scalar parameter cannot correct opposing miscalibration directions. These findings have direct implications for spectroscopic follow-up prioritization under the alert-to-spectrum mismatch of the Rubin era.
\end{abstract}

\keywords{Transient sources --- Sky surveys --- Astronomical methods --- Astrostatistics --- Classification --- Astronomy data analysis}

% ============================================================
\section{Introduction} \label{sec:intro}
% ============================================================

The Vera C. Rubin Observatory's Legacy Survey of Space and Time (LSST) is designed to generate on the order of $10^7$ alerts per night, with the first public alerts issued on 2026 February 24 \citep{ivezic2019}. With global spectroscopic follow-up capacity remaining fixed at roughly $10^5$ spectra per year, the community faces a substantial alert-to-spectrum mismatch. The Zwicky Transient Facility \citep[ZTF;][]{bellm2019}, currently serving as the primary LSST precursor survey, has driven the development of automated alert brokers to triage this data deluge and identify the most scientifically valuable transients for follow-up.

Several broker systems now operate in production on the ZTF alert stream, including ALeRCE \citep{forster2021}, Fink \citep{moller2021}, ANTARES \citep{narayan2018}, and Lasair \citep{smith2019}. These systems ingest alerts in real time, extract features from light curves and contextual information, and deploy machine learning classifiers to assign class probabilities to each transient. Astronomers routinely use these probability outputs to make follow-up decisions: a transient classified as ``90\% SN~Ia'' is far more likely to receive telescope time than one classified at ``60\% SN~Ia.'' Additional frameworks such as NEEDLE \citep{sheng2024} extend classification to rare transient types including superluminous supernovae (SLSNe) and tidal disruption events (TDEs).

This decision-making process relies on an implicit but critical assumption: that the classifier's stated probabilities are \emph{calibrated}---that is, among all objects assigned a probability $p$ of belonging to a given class, a fraction $p$ actually belongs to that class \citep{dawid1982}. Miscalibrated probabilities lead directly to misallocated resources: an overconfident classifier wastes telescope time on objects it is wrong about, while an underconfident one causes astronomers to overlook genuinely interesting transients. Despite the operational importance of calibration, no prior work has evaluated calibration across multiple production brokers using standard metrics on real survey data.

The machine learning calibration literature has established that modern neural networks are systematically overconfident \citep{guo2017}, a finding that has driven the development of post-hoc calibration methods including temperature scaling \citep{guo2017}, Platt scaling \citep{platt1999}, and Dirichlet calibration \citep{kull2019}. However, Random Forests---the architecture underlying several production broker classifiers---exhibit the opposite calibration behavior: ensemble averaging over bagged trees pushes predictions toward intermediate values, producing systematic underconfidence \citep{niculescumizil2005, bostrom2008}. Class imbalance corrections such as those used in Balanced Random Forests further distort probability estimates by inflating minority-class probabilities \citep{vandengoorbergh2022}. These well-characterized calibration properties of different architectures have direct relevance for astronomical classifiers but have not been evaluated in this domain.

Two recent works include limited calibration analysis. \citet{desoto2024} present per-class calibration curves for Superphot+ deployed on the ANTARES broker, noting overconfident SN~Ib/c and underconfident SN~Ia predictions but without computing formal calibration metrics. \citet{tung2025} evaluate calibration of ATCAT on simulated ELAsTiCC data, applying post-hoc corrections. However, no prior work has evaluated calibration across multiple production brokers using standard metrics on real survey data.

In this paper, we present the first such audit. We evaluate four classifiers across three production systems: ALeRCE's Balanced Random Forest light curve classifier \citep[lc\_classifier v1.1.13;][]{sanchezsaez2021}, Fink's Random Forest and SuperNNova classifiers \citep{leoni2022, moller2020}, and the NEEDLE hybrid CNN+DNN framework \citep{sheng2024}. We restrict our analysis to classifiers that output explicit class probabilities; other brokers such as Lasair, ANTARES, and AMPEL employ crossmatching or feature-based approaches that do not produce probability estimates amenable to calibration analysis. Using spectroscopically confirmed transients from the ZTF Bright Transient Survey \citep[BTS;][]{fremling2020, perley2020} as ground truth, we compute ECE with equal-mass binning \citep{roelofs2022}, classwise reliability diagrams \citep{nixon2019}, and Brier scores, then apply temperature scaling with five-fold cross-validation. We find that all four classifiers exhibit significant miscalibration, and that temperature scaling is effective for ALeRCE but inappropriate or counterproductive for the other systems.

The paper is organized as follows. Section~\ref{sec:background} reviews relevant background on calibration metrics and post-hoc methods. Section~\ref{sec:data} describes our data sources and sample construction. Section~\ref{sec:methods} presents our calibration analysis methodology. Section~\ref{sec:results} reports our findings for each classifier. Section~\ref{sec:discussion} interprets the results in the context of the calibration literature and discusses implications for LSST. Section~\ref{sec:conclusions} summarizes our conclusions.


% ============================================================
\section{Background} \label{sec:background}
% ============================================================

\subsection{Calibration and its Importance}

A probabilistic classifier is said to be \emph{calibrated} if, for all predicted probability values $p$, the empirical frequency of the positive class among predictions with confidence $p$ equals $p$ \citep{dawid1982}. Formally, for a classifier $f$ producing predictions $\hat{Y}$ with associated confidence $\hat{P}$:
\begin{equation}
P(\hat{Y} = Y \mid \hat{P} = p) = p \quad \forall p \in [0,1].
\end{equation}
Calibration is distinct from discrimination (the ability to separate classes). A classifier can achieve high accuracy while being poorly calibrated, and vice versa.

\subsection{Calibration Metrics}

The \emph{Expected Calibration Error} \citep[ECE;][]{naeini2015} is the most widely used scalar summary of calibration quality. It partitions predictions into $M$ bins based on confidence and computes the weighted average of the absolute difference between accuracy and confidence in each bin:
\begin{equation}
\text{ECE} = \sum_{m=1}^{M} \frac{|B_m|}{n} \left| \text{acc}(B_m) - \text{conf}(B_m) \right|,
\end{equation}
where $B_m$ is the set of samples in bin $m$, $\text{acc}(B_m)$ is the empirical accuracy, and $\text{conf}(B_m)$ is the mean confidence. We use $M = 15$ bins following \citet{guo2017}.

The choice of binning strategy significantly affects ECE estimates. Equal-width binning places bin boundaries at uniform intervals along the confidence axis, which can leave many bins empty or near-empty for classifiers with skewed confidence distributions. \citet{roelofs2022} demonstrated that equal-mass (quantile) binning, which places approximately the same number of samples in each bin, produces lower-bias ECE estimates and recommended against equal-width binning. We adopt equal-mass binning throughout.

Standard ECE considers only the top-predicted class and can mask miscalibration of individual classes. \citet{nixon2019} introduced the \emph{Static Calibration Error} (SCE) and \emph{Adaptive Calibration Error} (ACE), which evaluate calibration of all $K$ class probabilities independently. We report classwise ECE following this approach.

The \emph{Brier score} \citep{brier1950} provides a complementary proper scoring rule:
\begin{equation}
\text{BS} = \frac{1}{n} \sum_{i=1}^{n} \sum_{k=1}^{K} (p_{ik} - y_{ik})^2,
\end{equation}
where $p_{ik}$ is the predicted probability for class $k$ and $y_{ik}$ is the indicator for the true class.

\subsection{Temperature Scaling}

Temperature scaling \citep{guo2017} is a single-parameter post-hoc calibration method that rescales the logit vector $\mathbf{z}$ by a learned temperature $T > 0$:
\begin{equation}
\hat{q}_k = \frac{\exp(z_k / T)}{\sum_{j} \exp(z_j / T)}.
\end{equation}
When $T > 1$, the output distribution is softened (reducing overconfidence); when $T < 1$, it is sharpened (reducing underconfidence). The temperature is fitted by minimizing the negative log-likelihood (NLL) on a held-out calibration set, as NLL is a strictly proper scoring rule \citep{gneiting2007}.

Critically, temperature scaling applies a single global transformation and therefore cannot correct class-asymmetric miscalibration---where some classes are overconfident while others are underconfident. \citet{kull2019} demonstrated that temperature scaling can appear to produce well-calibrated confidence-level reliability diagrams while individual classes remain miscalibrated, and proposed Dirichlet calibration as a more flexible alternative.

\subsection{Calibration of Random Forests}

\citet{niculescumizil2005} showed that Random Forests produce systematically underconfident predictions due to the averaging of base learner probabilities: if the true probability is near 0 or 1, the ensemble average is biased toward 0.5 because only consensus among all trees can produce extreme probabilities. This effect is amplified by feature subsetting, which increases base learner variance. \citet{bostrom2008} confirmed this finding and proposed sigmoid calibration to correct it. \citet{loecher2024} demonstrated through simulation that Random Forest probability estimates ``are always underconfident'' across 192 scenarios.

Balanced Random Forests \citep{chen2004} introduce an additional calibration-distorting mechanism. By undersampling the majority class during training, they alter the effective class prior, which \citet{vandengoorbergh2022} showed produces ``strong and systematic overestimation of the probability for the minority class.'' The combined effect of standard RF underconfidence and minority-class inflation from balancing creates a complex but predictable calibration profile.

\subsection{Effects of Class Weighting on Calibration}

Many astronomical classifiers employ inverse-frequency class weighting or similar imbalance correction strategies to improve minority-class recall. \citet{vandengoorbergh2022} demonstrated that all common imbalance corrections---including random undersampling, oversampling, and SMOTE---harm probability calibration. \citet{fernando2021} confirmed that ``probability estimates of minority classes are unreliable'' when class weighting is applied. This is directly relevant to NEEDLE, which employs inverse-frequency weighting with extreme ratios ($\sim$80:1 for common supernovae versus TDEs).


% ============================================================
\section{Data} \label{sec:data}
% ============================================================

\subsection{Spectroscopic Ground Truth}

We use the ZTF Bright Transient Survey \citep[BTS;][]{fremling2020, perley2020} as our primary source of spectroscopic ground truth. BTS provides systematic spectroscopic classification of transients detected by ZTF, yielding 7,167 objects with spectroscopic type assignments. We map the $\sim$30 BTS spectroscopic subtypes to the ALeRCE taxonomy using a conservative mapping: SN~Ia and its subtypes (Ia-91T, Ia-91bg, Ia-02cx, Ia-CSM) map to SNIa; SN~Ib, Ic, Ic-BL, Ib/c, and Ibn map to SNIbc; SN~II, IIn, IIb, and IIP map to SNII; SLSN-I and SLSN-II map to SLSN; and TDE maps to TDE. Objects with types not present in any classifier taxonomy (novae, LRN, LBV) are excluded, yielding 7,116 mapped objects.

\subsection{Stratified Sampling}

Because rare classes (SLSN, TDE) constitute $<$2\% of the BTS catalog, we construct a stratified sample to ensure adequate representation for per-class calibration analysis. We cap the dominant class (SN~Ia) at 600 objects and include all available objects for rarer classes. This yields a sample of 1,436 objects with the distribution shown in Table~\ref{tab:sample}. Stratified sampling is appropriate for calibration analysis because ECE is a conditional measurement---it evaluates accuracy within confidence bins---and does not require the evaluation set to match the population class distribution \citep{guo2017}.

\begin{deluxetable}{lcc}
\tablecaption{Stratified Sample Composition \label{tab:sample}}
\tablehead{
\colhead{Class} & \colhead{Count} & \colhead{Fraction}
}
\startdata
SN~Ia & 600 & 41.8\% \\
SN~II & 341 & 23.7\% \\
SN~Ibc & 225 & 15.7\% \\
SLSN & 97 & 6.8\% \\
TDE & 39 & 2.7\% \\
Other & 134 & 9.3\%
\enddata
\end{deluxetable}

\subsection{Classifier Predictions}

\textbf{ALeRCE.} We query the ALeRCE API\footnote{\url{https://api.alerce.online}} for light curve classifier (lc\_classifier v1.1.13) probabilities. This Balanced Random Forest outputs a 15-class probability vector spanning transient and variable star categories \citep{sanchezsaez2021}. The 2025 TDE expansion described in \citet{pavezherrera2025} postdates our analysis; we evaluate the pre-TDE lc\_classifier v1.1.13. We successfully retrieve predictions for 1,149 of 1,436 sample objects. Since ALeRCE does not include a TDE class in v1.1.13, we exclude 35 TDE objects, yielding 1,114 objects for analysis. We restrict to the four transient classes (SN~Ia, SN~Ibc, SN~II, SLSN) and renormalize the probability vectors to sum to unity over these classes.

\textbf{Fink.} We query the Fink API\footnote{\url{https://api.fink-portal.org}} for two binary classifiers: the Random Forest SN~Ia classifier (\texttt{rf\_snia\_vs\_nonia}) and SuperNNova \citep[\texttt{snn\_snia\_vs\_nonia};][]{moller2020}. Both output a single score $\in [0, 1]$ representing $P(\text{SN~Ia})$. We retrieve the most recent alert for each object (corresponding to the most complete light curve), obtaining predictions for 1,237 objects.

\textbf{NEEDLE.} Unlike ALeRCE and Fink, NEEDLE \citep{sheng2024} does not provide a public API. We perform local inference using the five pre-trained models from the \texttt{lasair\_th\_r\_model} family, applied to the publicly available dataset. Each model was trained with a different random held-out test set recorded in \texttt{testset\_obj.json}. We map test set ZTF IDs to dataset positions via the provided hash table, yielding 429 predictions across 278 unique objects in three classes (SN, SLSN-I, TDE).


% ============================================================
\section{Methods} \label{sec:methods}
% ============================================================

\subsection{Calibration Evaluation}

For each classifier, we compute:
\begin{enumerate}
    \item ECE with equal-mass binning (15 bins) and, for comparison, equal-width binning.
    \item Classwise ECE following \citet{nixon2019}, evaluating calibration of each class probability independently.
    \item Reliability diagrams showing empirical accuracy versus mean confidence per bin.
    \item Brier score as a proper scoring rule.
    \item Bootstrap confidence intervals (1,000 replicates, 95\% level) on ECE.
\end{enumerate}

For binary classifiers (Fink RF and SuperNNova), confidence is defined as $\max(p, 1-p)$ and predictions are thresholded at $p = 0.5$.

For Fink classifiers, which implement implicit selective classification through prediction-eligibility gates \citep{geifman2017}, we report calibration metrics both \emph{unconditionally} (computed on all API outputs including score $= 0$) and \emph{conditionally} (restricted to score $> 0$ cases where predictions were made). The unconditional metric reflects the user experience when zeros are naively interpreted as probabilities; the conditional metric isolates actual classifier calibration on the predicted subset. This distinction is essential for systems where abstention is conflated with low-confidence prediction in the API output format.

\subsection{Temperature Scaling}

We fit the temperature parameter $T$ by minimizing negative log-likelihood using bounded scalar optimization ($T \in [0.01, 10.0]$). To avoid fitting and evaluating on the same data, we use five-fold stratified cross-validation: $T$ is fitted on four folds and ECE is evaluated on the held-out fold, with results averaged across folds. Stratified splitting ensures rare classes (SLSN, 6.8\% of sample) are represented in all folds.

For Random Forest classifiers, which do not natively produce logits, we convert vote-fraction probabilities to pseudo-logits via $z_k = \log p_k$ before applying temperature scaling. We note that this conversion assumes a log-linear relationship that may not hold precisely for RF outputs.

We apply three decision rules to determine whether temperature scaling should be recommended:
\begin{enumerate}
    \item \textbf{Bound detection:} If $T$ reaches the optimizer bound ($T \geq 9.5$ or $T \leq 0.02$), scaling is rejected as inappropriate.
    \item \textbf{Worsening detection:} If post-scaling ECE exceeds pre-scaling ECE, scaling is rejected.
    \item \textbf{Redundancy detection:} If pre-scaling ECE $< 0.1$, scaling is deemed unnecessary.
\end{enumerate}

For classifiers where global temperature scaling fails, we additionally evaluate per-class temperature scaling \citep{fang2020}, which fits an independent $T_k$ for each class and re-normalizes via softmax.


% ============================================================
\section{Results} \label{sec:results}
% ============================================================

Table~\ref{tab:summary} summarizes the calibration metrics for all four classifiers. We describe each in detail below.

\begin{deluxetable*}{llccccc}
\tablecaption{Calibration Audit Summary \label{tab:summary}}
\tablehead{
\colhead{Classifier} & \colhead{Subset} & \colhead{Architecture} & \colhead{$N$} & \colhead{ECE [95\% CI]} & \colhead{Post-calibration ECE} & \colhead{Coverage}
}
\startdata
ALeRCE & --- & BRF (15 classes) & 1114 & $0.271$ $[0.249, 0.296]$ & $0.097$ (T-scaling) & 100\% \\
\hline
\multirow{2}{*}{Fink RF} 
    & Raw API & \multirow{2}{*}{RF (2 classes)} & \multirow{2}{*}{1237} & $0.411$ $[0.382, 0.437]$ & N/A (abstention) & 100\% \\
    & Conditional ($\pi=1$) & & & $0.263$ $[0.211, 0.409]$ & $0.254$ (T-scaling) & 6.1\% \\
\hline
\multirow{2}{*}{Fink SNN} 
    & Raw API & \multirow{2}{*}{RNN (2 classes)} & \multirow{2}{*}{1237} & $0.198$ $[0.173, 0.226]$ & N/A (abstention) & 100\% \\
    & Conditional ($\pi=1$) & & & $0.183$ $[0.155, 0.220]$ & $0.118$ (T-scaling) & 64.1\% \\
\hline
NEEDLE & --- & CNN+DNN (3 classes) & 429 & $0.073$ $[0.063, 0.117]$ & $0.171$ (T worsens) & 100\% \\
\enddata
\tablecomments{ECE computed with equal-mass binning (15 bins). 95\% confidence intervals via bootstrap (1,000 replicates). For ALeRCE, post-calibration ECE is evaluated on held-out folds using stratified five-fold cross-validation. For Fink classifiers, ``Raw API'' treats all outputs including score $= 0$ as probabilities; ``Conditional'' restricts to score $> 0$ (predictions actually made). Coverage indicates the fraction of objects for which predictions were generated.}
\end{deluxetable*}

\subsection{ALeRCE: Systematic Underconfidence} \label{sec:alerce}

The ALeRCE light curve classifier exhibits systematic underconfidence across all transient classes (Figure~\ref{fig:alerce_reliability}). The overall accuracy is 0.742, but the mean confidence is only 0.471---a gap of $+0.271$ indicating that the classifier is substantially more accurate than its confidence scores suggest. The reliability diagram shows all bins lying above the diagonal, with confidence values compressed into the range $[0.3, 0.75]$ and never exceeding 0.75.

This underconfidence is consistent with the known calibration behavior of Random Forests \citep{niculescumizil2005, bostrom2008}. The Balanced Random Forest architecture amplifies this effect: by undersampling the majority class during training, it alters the effective prior, and the averaging of base learner predictions across a 15-class taxonomy further pushes probabilities toward uniform values.

Per-class analysis (Figure~\ref{fig:alerce_perclass}) reveals that underconfidence is pervasive: SN~Ia (ECE $= 0.195$, gap $= +0.335$), SN~Ibc (ECE $= 0.154$, gap $= +0.201$), and SN~II (ECE $= 0.150$, gap $= +0.393$) all exhibit strong underconfidence. SLSN is a notable exception with a slight overconfidence (gap $= -0.055$), potentially reflecting the Balanced Random Forest's inflation of minority-class probabilities \citep{vandengoorbergh2022}.

Temperature scaling with $T = 0.357 \pm 0.011$ (stratified five-fold CV) reduces ECE from 0.271 to 0.097---a 65\% improvement (Figure~\ref{fig:alerce_tempscaling}). The sharpening effect ($T < 1$) is the correct direction for underconfident classifiers, spreading the compressed confidence range $[0.3, 0.75]$ out to $[0.35, 1.0]$. We note that $T = 0.357$ is more aggressive than values typically reported in the neural network calibration literature (where $T > 1$ is standard for overconfident models), reflecting the severity of RF underconfidence when distributed across 15 classes.

The stratified cross-validation was essential: unstratified splitting underrepresents SLSN (6.8\% of sample), producing an artificially optimistic post-calibration ECE of 0.030. The honest stratified estimate of 0.097 accounts for calibration across all classes including rare ones.

\textbf{Operational impact.} Applying the recalibrated 0.8 probability threshold increases the number of usable SN~Ia candidates from 20 to 425 objects while maintaining 91\% precision, a 21-fold increase in follow-up targets. This demonstrates that miscalibration collapses the usable confidence axis along which astronomers allocate telescope time.

\begin{figure}
\plotone{fig1_alerce_reliability.pdf}
\caption{Reliability diagram for the ALeRCE light curve classifier showing systematic underconfidence: all bins lie above the perfect calibration diagonal. Confidence values are compressed into the range $[0.3, 0.75]$. ECE $= 0.271$ with equal-mass binning (15 bins). \label{fig:alerce_reliability}}
\end{figure}

\begin{figure*}
\plotone{fig2_alerce_perclass.pdf}
\caption{Per-class reliability diagrams for the ALeRCE classifier. SN~Ia, SN~Ibc, and SN~II show consistent underconfidence (bars above diagonal). SLSN shows slight overconfidence, potentially due to minority-class probability inflation from the Balanced Random Forest architecture. \label{fig:alerce_perclass}}
\end{figure*}

\begin{figure*}
\plotone{fig3_alerce_tempscaling.pdf}
\caption{ALeRCE reliability diagrams before (left) and after (right) temperature scaling with $T = 0.357$. ECE is reduced from 0.271 to 0.097 (65\% improvement) using stratified five-fold cross-validation. The sharpening effect spreads the compressed confidence distribution to span the full $[0.35, 1.0]$ range. \label{fig:alerce_tempscaling}}
\end{figure*}


\subsection{Fink: Implicit Selective Classification} \label{sec:fink}

Fink's two binary classifiers---a Random Forest \citep{leoni2022} and a recurrent SuperNNova network \citep{moller2020}---both exhibit implicit selective classification through prediction-eligibility mechanisms that determine whether the model produces a probabilistic output. When eligibility criteria are not met, the Fink API returns \texttt{score = 0.0}, which does not represent $P(\text{SN~Ia}) = 0$ but rather ``no prediction made.'' This design reflects a principled engineering choice: abstain rather than predict on incomplete or low-quality data. However, it creates a calibration-evaluation challenge because the raw API output conflates predictions with abstentions, leading to aggregate metrics that do not isolate classifier calibration from prediction coverage.

To address this, we evaluate Fink's classifiers using both \emph{unconditional} calibration metrics (treating all API outputs including \texttt{score = 0} as probability estimates) and \emph{conditional} metrics (restricting to cases where \texttt{score > 0}, i.e., the classifier actually ran). The unconditional metric reflects the user experience when zeros are naively interpreted as probabilities; the conditional metric isolates the calibration of the classifier itself. This distinction is essential for understanding whether observed miscalibration arises from the prediction model or from prediction-coverage design choices.

\subsubsection{Prediction Eligibility and Abstention Mechanisms}

The Random Forest classifier implements a documented eligibility gate: objects must have at least three detections per photometric filter to receive a prediction \citep{leoni2022}. When this criterion is not met, the API returns \texttt{score = 0.0}. Figure~\ref{fig:fink_rf_histogram} shows the score distribution on our 1,237-object sample: 1,161 objects (93.9\%) return exactly zero, with a gap of 0.031 between zero and the minimum non-zero score. This gap signature is diagnostic of a hard threshold rather than a continuous low-probability tail, confirming that RF zeros represent abstentions rather than genuine $P(\text{SN~Ia}) \approx 0$ estimates.

SuperNNova's abstention mechanism is less documented but exhibits similar structure. Figure~\ref{fig:fink_snn_histogram} shows 444 objects (35.9\%) return exactly zero, with a gap of 0.005 to the first non-zero score and only 4 objects in the range $(0, 0.01)$. This suggests a possible eligibility threshold, though the smaller gap and higher prediction coverage (64.1\% vs.\ RF's 6.1\%) indicate a less conservative gate. The mechanism may involve minimum light curve length, signal-to-noise requirements, or model confidence thresholds, but without access to Fink's internal implementation we classify this as ``likely abstention'' pending verification.

Mathematically, both classifiers implement a form of \emph{selective classification} \citep{geifman2017}: the system computes a coverage function $\pi(x) \in \{0, 1\}$ indicating prediction eligibility, where $\pi(x) = 0$ triggers abstention. Unlike learning-to-defer frameworks \citep{mozannar2020,madras2018} where deferral policies are optimized jointly with prediction accuracy, Fink's abstention rules appear to be deterministic functions of data quality rather than learned confidence estimates.

\subsubsection{Unconditional Calibration (Raw API Experience)}

We first evaluate calibration as a naive user would experience it, treating all API outputs as probability estimates regardless of whether they represent predictions or abstentions. On our sample of $N = 1{,}237$ objects (527 SN~Ia, 710 not SN~Ia), the Random Forest returns zero for 1,161 cases, yielding an unconditional ECE of $0.411$ [95\% CI: $0.382$, $0.437$]. SuperNNova returns zero for 444 cases, with unconditional ECE of $0.198$ [95\% CI: $0.173$, $0.226$].

Figure~\ref{fig:fink_raw} shows reliability diagrams for both classifiers using all API outputs. The RF diagram is dominated by a single bin at confidence $\approx 1.0$ containing the 1,161 abstention cases treated as high-confidence predictions of ``not SN~Ia.'' This produces severe overconfidence: the bin shows 61\% accuracy at stated 100\% confidence. The SNN diagram similarly shows a concentration of zeros in the leftmost bin, but the effect is less extreme due to higher prediction coverage.

These unconditional ECE values describe the mismatch between stated confidence and accuracy when abstentions are treated as scores, a failure mode that Fink's documentation does not explicitly guard against. They do not, however, isolate the calibration properties of the prediction models themselves.

\subsubsection{Conditional Calibration (Predictions Only)}

Restricting to objects where the classifiers actually produced predictions (\texttt{score > 0}), we compute conditional ECE on the $\pi = 1$ subset. The Random Forest evaluated 76 objects (6.1\% coverage), achieving conditional ECE of $0.263$ [95\% CI: $0.211$, $0.409$] with 61.8\% accuracy at the 0.5 threshold. SuperNNova evaluated 793 objects (64.1\% coverage), achieving conditional ECE of $0.183$ [95\% CI: $0.155$, $0.220$] with 58.9\% accuracy.

Figure~\ref{fig:fink_conditional} shows reliability diagrams for the conditional subsets. The RF diagram reveals persistent miscalibration even after removing abstentions: bins in the $[0.6, 0.9]$ confidence range show systematic overconfidence, suggesting that the eligibility gate filters for data completeness but not prediction reliability. The small sample size (76 objects) yields wide confidence intervals, but the central estimate of ECE $= 0.263$ is substantially higher than the $< 0.05$ threshold typically considered well-calibrated.

SuperNNova's conditional diagram shows better-structured calibration with less severe overconfidence, though still above ideal. Notably, temperature scaling on the conditional subset yields substantial improvement: $T = 3.685$ reduces ECE from $0.183$ to $0.118$ (36\% improvement). In contrast, temperature scaling on RF's conditional subset provides minimal benefit ($0.263 \to 0.254$, 3\% improvement), suggesting that RF's miscalibration arises from limited discriminative power on the small eligible subset rather than a systematic temperature bias.

\subsubsection{Interpretation and Implications}

Fink's design implements a conservative eligibility-gating strategy that prioritizes precision over coverage: abstain when data quality is insufficient. This is conceptually aligned with selective classification frameworks \citep{geifman2017}, where systems learn when not to predict. However, three key differences distinguish Fink's approach from learning-to-defer paradigms \citep{mozannar2020}:

\begin{enumerate}
    \item \textbf{Hard-coded vs.\ learned deferral}: Fink's eligibility rules are deterministic thresholds (e.g., $\geq 3$ detections per filter) rather than learned functions of feature-space uncertainty.
    \item \textbf{Data quality vs.\ prediction confidence}: The RF gate filters on light curve completeness, not on model confidence or expected calibration error. Objects that barely meet the 3-epoch threshold receive predictions with ECE $= 0.263$, demonstrating that eligibility does not imply reliability.
    \item \textbf{Implicit vs.\ explicit abstention}: The API silently maps abstentions to \texttt{score = 0.0} without metadata indicating whether a zero represents ``confident not-SN-Ia'' or ``no prediction made.'' Users must reverse-engineer the abstention mechanism from score distributions.
\end{enumerate}

The operational consequence is a prediction-coverage tradeoff: RF achieves 6.1\% coverage with ECE $= 0.263$, while SNN achieves 64.1\% coverage with ECE $= 0.183$. Neither classifier is well-calibrated on its predicted subset in the raw form, though SNN becomes usable (ECE $= 0.118$) after temperature scaling. For astronomers building spectroscopic follow-up queues, this means that Fink scores---even when non-zero---require post-hoc recalibration before being used as decision thresholds.

The finding that conditional ECE differs substantially from unconditional ECE (RF: $0.263$ vs.\ $0.411$; SNN: $0.183$ vs.\ $0.198$) demonstrates that standard calibration metrics must account for prediction coverage in systems that implement abstention. Future work extending Fink's implicit selective classification to explicit learning-to-defer policies---where deferral decisions are optimized jointly with prediction accuracy and informed by expected spectroscopic utility---is a natural direction explored in forthcoming work on deferral-aware follow-up prioritization.

\begin{figure}
\plotone{diagnostic_random_forest_histogram.pdf}
\caption{Score distribution for Fink Random Forest on the BTS validation sample. The concentration of 1,161 objects (93.9\%) at exactly zero, with a gap of 0.031 to the minimum non-zero score, is diagnostic of a hard eligibility threshold rather than a continuous probability distribution. This confirms that \texttt{score = 0} represents abstention (``no prediction made'') rather than $P(\text{SN~Ia}) = 0$. \label{fig:fink_rf_histogram}}
\end{figure}

\begin{figure}
\plotone{diagnostic_supernnova_histogram.pdf}
\caption{Score distribution for Fink SuperNNova. The concentration of 444 objects (35.9\%) at zero with a gap of 0.005 to the first non-zero score suggests a likely abstention mechanism, though the smaller gap and higher prediction coverage (64.1\%) indicate a less conservative eligibility gate than the Random Forest. \label{fig:fink_snn_histogram}}
\end{figure}

\begin{figure}
\plotone{fig_fink_random_forest_raw.pdf}
\caption{Reliability diagram for Fink Random Forest (raw API, unconditional). The degenerate output distribution (94\% of scores at exactly 0.0) produces a misleading reliability diagram concentrated at confidence $\approx 1.0$. ECE $= 0.411$. \label{fig:fink_rf_raw}}
\end{figure}

\begin{figure}
\plotone{fig_fink_random_forest_conditional.pdf}
\caption{Reliability diagram for Fink Random Forest conditional on predictions (\texttt{score > 0}, $n = 76$ objects). Even after removing abstentions, the classifier exhibits ECE $= 0.263$, with systematic overconfidence in the $[0.6, 0.9]$ confidence range. This demonstrates that the eligibility gate filters for data completeness but not prediction reliability. \label{fig:fink_rf_conditional}}
\end{figure}

\begin{figure}
\plotone{fig_fink_supernnova_raw.pdf}
\caption{Reliability diagram for Fink SuperNNova (raw API, unconditional). Bars consistently below the diagonal indicate systematic overconfidence. ECE $= 0.198$. \label{fig:fink_snn_raw}}
\end{figure}

\begin{figure}
\plotone{fig_fink_supernnova_conditional.pdf}
\caption{Reliability diagram for Fink SuperNNova conditional on predictions (\texttt{score > 0}, $n = 793$ objects). The classifier achieves ECE $= 0.183$ on the predicted subset, which temperature scaling ($T = 3.685$) improves to 0.118, making it usable for spectroscopic follow-up prioritization after post-hoc recalibration. \label{fig:fink_snn_conditional}}
\end{figure}


\subsection{NEEDLE: Hidden Class-Asymmetric Miscalibration} \label{sec:needle}

NEEDLE's aggregate calibration appears acceptable: ECE $= 0.073$, with the reliability diagram showing bars that generally track the diagonal (Figure~\ref{fig:needle_aggregate}). However, per-class analysis reveals a qualitatively different picture.

Figure~\ref{fig:needle_perclass} shows the classwise reliability diagrams. SN predictions are underconfident (accuracy 0.929 versus mean confidence 0.870, gap $= +0.059$), while SLSN-I predictions are overconfident (accuracy 0.690 versus mean confidence 0.846, gap $= -0.156$). TDE predictions show mild overconfidence (accuracy 0.802 versus mean confidence 0.823, gap $= -0.021$). The aggregate ECE conceals this asymmetry because the underconfident SN class (which dominates by sample size) partially cancels the overconfident SLSN-I class in the weighted average.

This class-asymmetric miscalibration is a direct consequence of NEEDLE's training procedure. \citet{sheng2024} employ inverse-frequency class weighting to address the extreme imbalance between common supernovae and rare transients. With approximately 80$\times$ more SN training examples than TDE examples, the effective weight for rare classes is $\sim$80$\times$ higher per sample. \citet{vandengoorbergh2022} demonstrated that such corrections systematically inflate minority-class probability estimates, which is precisely the pattern we observe.

Global temperature scaling with $T = 1.551$ \emph{worsens} ECE from 0.073 to 0.171. This failure is expected: a single $T$ cannot simultaneously soften overconfident SLSN-I predictions ($T > 1$ needed) and sharpen underconfident SN predictions ($T < 1$ needed). \citet{kull2019} documented this fundamental limitation, noting that temperature scaling ``cannot learn to act differently on different classes.'' Per-class temperature scaling yields $T_{\text{SN}} = 1.85$, $T_{\text{SLSN-I}} = 2.03$, and $T_{\text{TDE}} = 0.54$, confirming the opposing calibration directions, but does not improve aggregate ECE due to the small sample sizes involved (75 SLSN-I, 105 TDE predictions across 5 models).

Silent failure analysis reveals 17 predictions at $\geq$90\% confidence that are incorrect, concentrated on three SLSN-I objects (ZTF20aagikvv, ZTF19abzoyeg, ZTF21aakjkec) that are consistently misclassified as SN with near-100\% confidence across all five models. This failure pattern is consistent with the class-weighting mechanism: the rare-class inflation boosts SLSN-I probabilities during training, but the effect is strongest on ambiguous cases near the decision boundary. The result is a calibration failure that concentrates precisely where follow-up resources are scarcest and misclassification is most expensive.

\begin{figure}
\plotone{fig_needle_reliability.pdf}
\caption{Aggregate reliability diagram for NEEDLE. ECE $= 0.073$ suggests acceptable calibration, but this conceals class-level miscalibration (see Figure~\ref{fig:needle_perclass}). \label{fig:needle_aggregate}}
\end{figure}

\begin{figure*}
\plotone{fig_needle_perclass.pdf}
\caption{Per-class reliability diagrams for NEEDLE revealing class-asymmetric miscalibration hidden by the aggregate ECE. SN is underconfident (bars above diagonal); SLSN-I is overconfident (high-confidence bars below diagonal, especially near 1.0); TDE is mildly overconfident. This asymmetry is driven by inverse-frequency class weighting during training. \label{fig:needle_perclass}}
\end{figure*}

\subsection{Comparative Overview}

Figure~\ref{fig:summary} presents the ECE comparison across all four classifiers. None achieves the ``well-calibrated'' threshold of ECE $< 0.05$ in their raw form. Only ALeRCE's miscalibration is addressable by temperature scaling, reducing its ECE below the ``acceptable'' threshold of ECE $< 0.10$. Fink's SuperNNova becomes usable (ECE $= 0.118$) after temperature scaling on the conditional subset. The remaining classifiers require either architectural changes (Fink RF) or class-aware calibration approaches (NEEDLE).

\begin{figure*}
\plotone{fig8_comparison.pdf}
\caption{ECE comparison across all four production classifiers. None meets the ``well-calibrated'' threshold (ECE $< 0.05$, dotted line) in raw form. Only NEEDLE falls below the ``acceptable'' threshold (ECE $< 0.10$, dashed line) in aggregate, but this masks per-class miscalibration. Error bars show 95\% bootstrap confidence intervals. \label{fig:summary}}
\end{figure*}


% ============================================================
\section{Discussion} \label{sec:discussion}
% ============================================================

\subsection{Architecture Determines Calibration Failure Mode}

Our results demonstrate a clear relationship between classifier architecture and the nature of miscalibration. The ALeRCE Balanced Random Forest exhibits underconfidence consistent with the well-established calibration properties of ensemble methods \citep{niculescumizil2005, bostrom2008, loecher2024}. This is distinct from the overconfidence typically found in deep neural networks \citep{guo2017}, highlighting the importance of architecture-specific calibration analysis rather than assuming universal overconfidence.

NEEDLE's class-asymmetric miscalibration represents a qualitatively different failure mode that is invisible to standard aggregate ECE. This finding aligns with \citet{nixon2019}, who demonstrated that aggregate calibration metrics can conceal per-class miscalibration, and with \citet{kull2019}, who showed that confidence-level reliability diagrams can appear well-calibrated even when classwise diagrams reveal significant problems. We recommend that future calibration analyses of multi-class astronomical classifiers always include per-class evaluation.

\subsection{The Failure of Temperature Scaling as a Diagnostic}

The failure of global temperature scaling for NEEDLE and Fink is itself informative. When temperature scaling worsens calibration, it indicates that the miscalibration is structural---arising from the model's architecture or training procedure---rather than a simple uniform shift in confidence. For NEEDLE, the opposing per-class temperatures ($T_{\text{SN}} = 1.85$ versus $T_{\text{TDE}} = 0.54$) directly trace the effect of inverse-frequency class weighting, confirming the mechanism identified by \citet{vandengoorbergh2022}. For Fink SuperNNova, the bound-hitting behavior ($T = 10$) on the unconditional dataset reveals that the bimodal score distribution (36\% exact zeros plus a continuous component) is fundamentally incompatible with temperature scaling's assumption of a smooth logit space.

These diagnostic patterns suggest a practical recommendation: apply temperature scaling not only as a correction method but as a probe of miscalibration structure. If $T \approx 1$, the classifier is approximately calibrated. If $T$ converges to a finite value away from 1 and ECE improves, the miscalibration is uniform and fixable. If $T$ hits a bound or worsens ECE, deeper architectural or training-level interventions are needed.

\subsection{Selective Classification and Implicit Deferral}

Fink's eligibility gating represents an early example of selective classification in astronomical ML pipelines. The system effectively decides when \emph{not} to predict, though the deferral rule is deterministic (data completeness threshold) rather than probabilistic. This finding connects directly to learning-to-defer frameworks \citep{mozannar2020,madras2018}, which optimize deferral policies jointly with prediction accuracy to balance coverage against error rates under resource constraints.

Making Fink's implicit abstention explicit and learnable---for instance, by training a model to predict expected calibration error or spectroscopic utility as a function of light curve features---is a natural extension. Such a framework would replace hard-coded data-quality thresholds with learned confidence estimates that account for both model uncertainty and the value of spectroscopic confirmation. We explore this direction in forthcoming work on deferral-aware spectroscopic follow-up prioritization, where the goal is not merely to classify transients but to optimize scientific utility per telescope hour under the alert-to-spectrum mismatch of the Rubin era.

\subsection{Implications for Spectroscopic Follow-up}

Uncalibrated probabilities have direct consequences for follow-up allocation. An astronomer using ALeRCE's raw probabilities and requiring $p > 0.8$ for follow-up would observe very few objects---the classifier's confidence never exceeds 0.75---despite 74\% of its predictions being correct. After temperature scaling, the calibrated confidence distribution spans $[0.35, 1.0]$, enabling meaningful probability-based thresholding. The 21-fold increase in usable candidates (20 $\to$ 425 at 91\% precision) demonstrates that miscalibration collapses the confidence axis along which astronomers allocate telescope time.

For NEEDLE, the practical concern is SLSN-I overconfidence. When the classifier reports 95\% confidence in an SLSN-I classification, the actual accuracy is only 69\%. Given the scientific value of SLSNe and TDEs for probing extreme physics, overconfident classification of these rare objects can lead to wasted spectroscopic resources---precisely the failure mode that calibration is designed to prevent.

\subsection{Implications for LSST}

With the Rubin Observatory's public alert stream now operational, the calibration problems we identify become increasingly urgent. The greater diversity of transient types, the increased class imbalance (rarer classes become proportionally rarer in a larger survey), and the higher volume of alerts all exacerbate miscalibration risk. Our findings suggest three concrete recommendations:

First, broker teams should publish calibration metrics alongside classification performance. ECE, reliability diagrams, and Brier scores should be standard evaluation outputs, not afterthoughts.

Second, post-hoc calibration should be architecture-aware. Temperature scaling is appropriate for Random Forest classifiers but can be counterproductive for class-weighted neural networks. Class-aware methods such as Dirichlet calibration \citep{kull2019} or focal calibration \citep{mukhoti2020} may be more appropriate for classifiers with extreme class imbalance.

Third, calibration should be evaluated on real survey data, not only on simulated datasets. Our finding that Fink's RF produces 94\% zero scores when evaluated on BTS objects---behavior that would not appear in simulated data designed to match the classifier's training distribution---illustrates the gap between simulated and operational calibration.

\subsection{Robustness Checks}

We performed two robustness checks to verify that our findings are not artifacts of data processing choices. Script 05 (renormalization audit) confirmed that restricting ALeRCE's 15-class output to the 4-class transient subset does not drive the observed miscalibration---the renormalization step does not introduce calibration artifacts. Script 06 (duplicate audit) for NEEDLE found that deduplicating objects appearing in multiple test folds yields dedup ECE $= 0.048$ versus 0.073 for the full set. The class-asymmetric miscalibration pattern (SN underconfident, SLSN-I overconfident) persists in the deduplicated set with 100\% inter-model class agreement, confirming that the finding is robust to evaluation-set construction choices.

\subsection{Limitations}

Several limitations should be noted. Our ground truth is limited to spectroscopically confirmed objects from BTS, which is itself magnitude-limited and subject to selection effects. The stratified sampling may not capture the full range of classifier behavior across the alert stream. For NEEDLE, our evaluation uses predictions from five models on their respective held-out test sets, which may not represent the classifier's performance on the broader alert stream.

The small sample sizes for rare classes (97 SLSN and 39 TDE in the BTS sample; 75 SLSN-I and 105 TDE predictions for NEEDLE) limit the statistical power of per-class calibration analysis. More flexible calibration methods such as isotonic regression or Dirichlet calibration require larger samples to fit reliably \citep{vaicenavicius2019}, which is why we do not attempt them here.

Finally, our temperature scaling converts RF vote fractions to pseudo-logits via $z_k = \log p_k$, an assumption that may not be optimal for non-softmax-based classifiers. The unusually aggressive temperature ($T = 0.357$) may partly reflect this modeling choice rather than the true severity of miscalibration.


% ============================================================
\section{Conclusions} \label{sec:conclusions}
% ============================================================

We have presented the first systematic calibration audit of production machine learning classifiers deployed on ZTF transient alert streams. Our key findings are:

\begin{enumerate}
    \item No prior work has evaluated calibration across multiple production brokers using standard metrics on real survey data. We find significant miscalibration across all audited systems, with distinct failure modes tied to classifier architecture.

    \item ALeRCE's Balanced Random Forest is systematically underconfident (ECE $= 0.271$), consistent with known RF calibration behavior. Temperature scaling with $T = 0.357$ reduces ECE to 0.097 (65\% improvement) using stratified cross-validation. Applying a recalibrated 0.8 probability threshold increases usable SN~Ia candidates from 20 to 425 objects while maintaining 91\% precision, a 21-fold operational gain.

    \item Fink implements implicit selective classification: the Random Forest abstains on 94\% of objects via a minimum-epoch eligibility gate, while SuperNNova abstains on 36\%. Conditional calibration metrics (computed only on score $> 0$ subsets) reveal ECE $= 0.263$ for RF and ECE $= 0.183$ for SNN, both of which exceed well-calibrated thresholds. Unconditional metrics (treating abstentions as scores) yield ECE $= 0.411$ and 0.198 respectively, demonstrating that aggregate calibration metrics must account for prediction coverage in systems with abstention mechanisms. Temperature scaling on SNN's conditional subset achieves ECE $= 0.118$, making it usable for follow-up prioritization after recalibration.

    \item NEEDLE's aggregate ECE of 0.073 masks class-asymmetric miscalibration driven by inverse-frequency class weighting. SLSN-I predictions are overconfident (accuracy 69\% at mean confidence 85\%), while SN predictions are underconfident. Global temperature scaling worsens calibration (ECE $0.073 \to 0.171$) because it cannot correct opposing miscalibration directions simultaneously. Silent failures concentrate on boundary cases where class weighting inflates confidence but the classifier is least equipped to handle.

    \item The failure mode of temperature scaling is itself diagnostic: worsening ECE indicates structural miscalibration requiring architectural intervention, not post-hoc correction.
\end{enumerate}

As time-domain astronomy operates under the alert-to-spectrum mismatch of the Rubin era, calibrated probabilities are essential for optimal spectroscopic resource allocation. We recommend that broker teams adopt calibration evaluation as a standard component of classifier assessment, using proper scoring rules, per-class reliability analysis, and conditional evaluation for systems implementing abstention.

Our analysis code and results are publicly available at \url{https://github.com/pallavi-2000/transient-calibration-audit}.

\section*{Data and Code Availability}

All analysis code, data acquisition scripts, and results are publicly available at \url{https://github.com/pallavi-2000/transient-calibration-audit}. The repository includes scripts to reproduce all figures and tables from the raw data.

% ============================================================
\begin{acknowledgments}
We thank Matt Nicholl for encouragement and for suggesting connections with the NEEDLE team, and Xinyue Sheng for making the NEEDLE models and dataset publicly available. This work made use of the ALeRCE broker \citep{forster2021}, the Fink broker \citep{moller2021}, and data from the ZTF Bright Transient Survey \citep{fremling2020, perley2020}. This research has made use of NASA's Astrophysics Data System.
\end{acknowledgments}


% ============================================================
\software{
\texttt{NumPy} \citep{harris2020},
\texttt{SciPy} \citep{virtanen2020},
\texttt{pandas} \citep{mckinney2010},
\texttt{Matplotlib} \citep{hunter2007},
\texttt{TensorFlow} \citep{abadi2016}
}


% ============================================================
\bibliography{references2}
\bibliographystyle{aasjournal}

\end{document}